% Figure peps_dyadic_hv_routing.
% Owner: chapter/ch24_peps_dyadic_routing.tex (def:peps_dyadic_hv_routing).
% Used by: gapped-2d-notes PP Sec. 10.
% Source: arXiv OpenAI2026PolynomialPEPS, 07-assembly.tex, lines
% 101--160 (anchors eq. (anchors), proof of lem:routing); the source has
% no figure for this step.
% Ink-to-index: the drawing carries no tensor indices. Every dot is a
% point of the padded square {0,...,7}^2; cell (r,c) is the point
% (x,y)=(c-1,r-1), with y increasing downward on the page. Every point is
% a scale-zero anchor. The square marks the scale-3 anchor (4,4); the
% four rounded atoms mark the scale-2 anchors {2,6}^2; the blue dots mark
% the scale-1 anchors {1,3,5,7}^2.
% The directed string is one link from P=(2,2) (scale 2, block indices
% (0,0)) to Q=(7,5) (scale 1, block indices (3,2)): horizontal along the
% row y=2 of P to the column x=7 of Q, then vertical to Q. Its separation
% (5,3) is within c times either block side for c=3.
% The shaded enclosure is the genuine square {0,...,4}^2 (L=5). The fold
% f(x)=x for x<=4 and f(x)=8-x for x>=5 sends Q to f(Q)=(1,3), cell
% (4,2). The folded walk retraces itself (x: 2,3,4,3,2,1 on y=2, then
% y: 2,3,4,3 on x=1), so it is described here and not drawn.
% No physical or virtual legs; empty boundary signature.
\begin{tenkzequation}
    \tenkzkernel
    \begin{tenkz}[rows={ket,ket,ket,ket,ket,ket,ket,ket}, cols=8,
            bonds=none]
        \tndeclare{species}{scale-one}{hue=source:blue}
        \tndeclare{species}{route}{hue=source:green!50!black}
        \tn[at=(5,5), skin=box]{}
        \tn[at=(3,3), skin=ring, name=P, label pos=nw]{P}
        \tn[at=(3,7), skin=ring]{}
        \tn[at=(7,3), skin=ring]{}
        \tn[at=(7,7), skin=ring]{}
        \tn[at=(2,2), species=scale-one]{}
        \tn[at=(2,4), species=scale-one]{}
        \tn[at=(2,6), species=scale-one]{}
        \tn[at=(2,8), species=scale-one]{}
        \tn[at=(4,2), species=scale-one]{}
        \tn[at=(4,4), species=scale-one]{}
        \tn[at=(4,6), species=scale-one]{}
        \tn[at=(4,8), species=scale-one]{}
        \tn[at=(6,2), species=scale-one]{}
        \tn[at=(6,4), species=scale-one]{}
        \tn[at=(6,6), species=scale-one]{}
        \tn[at=(6,8), species=scale-one, name=Q, label pos=e]{Q}
        \tn[at=(8,2), species=scale-one]{}
        \tn[at=(8,4), species=scale-one]{}
        \tn[at=(8,6), species=scale-one]{}
        \tn[at=(8,8), species=scale-one]{}
        \tnwire[kind=string, species=route, route=orth, dir=to,
            name=link, via={(3,8)}]{P}{Q}
        \tnmark[form=label, label pos=s]{(4,2)}{$f(Q)$}
        \tnmark[form=enclosure, species=selected, tint, label pos=nw]
        {(1,1) .. (5,5)}{$L\times L$}
    \end{tenkz}
\end{tenkzequation}
